% Scientific Paper Writer — JAMA Network Template
% Usage: Place in manuscript/ directory, compile with pdflatex + bibtex
% JAMA Network Open accepts single-column submissions; two-column is for final typesetting.
\documentclass[12pt,a4paper]{article}

% Packages
\usepackage[margin=1in]{geometry}
\usepackage{setspace}\doublespacing
\usepackage{times}
\usepackage{graphicx}
\usepackage[labelfont=bf]{caption}
\usepackage{booktabs}
\usepackage{hyperref}
\usepackage[numbers,sort&compress]{natbib}
\usepackage{lineno}\linenumbers
\usepackage{amsmath}
\usepackage{multicol}
\usepackage{framed}
\usepackage{enumitem}

% Title page
\title{\textbf{TITLE}}
\author{
  Author One, MD\textsuperscript{1},
  Author Two, PhD\textsuperscript{2} \\\\
  \small\textsuperscript{1}Department, Institution, City, Country \\\\
  \small\textsuperscript{2}Department, Institution, City, Country \\\\
  \small Corresponding author: Name, MD, Department, Institution, Address. (email@example.com)
}
\date{}

\begin{document}

\maketitle

% ============================================================
% Key Points (required for JAMA — outside abstract word limit)
% ============================================================
\begin{framed}
\noindent\textbf{Key Points}

\textbf{Question:} What is the [clinical question]?

\textbf{Findings:} In this [study design] of [N participants/studies], [primary finding with effect estimate, 95\% CI, and p-value]. [Secondary finding if applicable].

\textbf{Meaning:} [1--2 sentences on clinical/policy implications].
\end{framed}

\vspace{1em}

% Word count and counts
\noindent\textbf{Word count:} XXXX words \\
\textbf{Number of tables:} X \\
\textbf{Number of figures:} X \\
\textbf{Number of references:} XX

\newpage

% ============================================================
% Structured Abstract (max 350 words for JAMA Network Open)
% ============================================================
\begin{abstract}
\noindent\textbf{Importance:} \\[0.5em]
\textbf{Objective:} To \ldots \\[0.5em]
\textbf{Data Sources:} \\[0.5em]
\textbf{Study Selection:} \\[0.5em]
\textbf{Data Extraction and Synthesis:} \\[0.5em]
\textbf{Main Outcomes and Measures:} \\[0.5em]
\textbf{Results:} \\[0.5em]
\textbf{Conclusions and Relevance:}
\end{abstract}

\newpage

% ============================================================
% Main text
% ============================================================
\section{Introduction}
% INSERT INTRODUCTION HERE
% JAMA guidelines: 2-3 paragraphs, ending with a clear statement of the study objective.

\section{Methods}
% INSERT METHODS HERE
% Include subsections as appropriate for study type.
% End with: "This study followed the [PRISMA/STROBE/CONSORT] reporting guideline."
% Ethics: "This study was approved by [IRB name] (approval No. XXX)."

\section{Results}
% INSERT RESULTS HERE

\section{Discussion}
% INSERT DISCUSSION HERE
% JAMA guidelines: 4-5 paragraphs.
% Paragraph 1: Summary of principal findings
% Paragraph 2-3: Comparison with prior literature
% Paragraph 4: Strengths and limitations
% Paragraph 5: Conclusions and implications

% ============================================================
% References — JAMA uses superscript numbered citations
% ============================================================
% In-text: use \textsuperscript{1} or \textsuperscript{1,2} or \textsuperscript{1-3}
\bibliographystyle{vancouver}
\bibliography{references}

% ============================================================
% Supplementary
% ============================================================
\newpage
\appendix
\section*{Supplement}

\subsection*{eMethods}
% INSERT SUPPLEMENTARY METHODS HERE

\subsection*{eTables}
% INSERT SUPPLEMENTARY TABLES HERE

\subsection*{eFigures}
% INSERT SUPPLEMENTARY FIGURES HERE

\end{document}
